\documentclass[conference]{IEEEtran}
\usepackage{amsmath, amssymb}
\usepackage{graphicx}
\usepackage{booktabs}
\usepackage{cite}

\title{Multi-Robot Infield Defense System Using Distributed Intercept and Relay Coordination}

\author{
\IEEEauthorblockN{Aman Khatri, Dena Shink, Jason Ballinger}
\IEEEauthorblockA{Indiana University Bloomington}
}

\begin{document}

\maketitle

\begin{abstract}
Multi-robot coordination systems operate in dynamic environments where robots must allocate tasks efficiently while maintaining system stability. This project proposes a distributed coordination framework inspired by baseball infield defense. Robots maintain formation, detect and intercept a moving ball, and determine optimal delivery strategies using decentralized decision-making. The system integrates trajectory prediction, interception optimization, and relay-based delivery to minimize total task completion time. The framework is implemented in a Python-based simulation environment to evaluate performance under varying conditions.
\end{abstract}

\section{Introduction}

Multi-robot systems (MRS) enable coordination, where robots must efficiently allocate tasks to solve complex tasks in dynamic environments such as warehouse automation, autonomous transportation, and search-and-rescue operations. A key challenge in these systems is dynamic task allocation, where robots must determine which agent should execute a task while maintaining coordination with the rest of the system.

Our project explores multi-robot coordination inspired by baseball infield defense. In this setting, multiple agents must quickly decide who should field a ground ball and how to move it toward a target base. This provides a structured scenario for studying decentralized coordination and role assignment.

We propose a distributed intercept–relay coordination framework in which robots can:
\begin{itemize}
    \item Maintain a defensive formation
    \item Detect and predict ball motion
    \item Independently compute interception costs
    \item Dynamically assign interception task without centralized control
\end{itemize}

\textbf{Updated Scope:} We are targeting a 2-Dimensional simulation of $N$ infield-style robots with kinematic motion, a ground ball modeled with constant planar acceleration, decentralized formation control on a ring communication graph, and time-to-intercept bidding for primary fielder selection. Our implementation uses Python with NumPy for dynamics and Matplotlib for visualization. A future, more complex implementation would ideally use a full physics engine rendering a 3-Dimensional simulation. However, we deemed this to be too complex given our time constraint on the project.

\textbf{Progress So Far:}
\begin{itemize}
    \item Fully developed the initial 2D simulation environment
    \item Fully developed the ball mechanisms, including motion, hitting, and foul ball randomness and logic.
    \item Partially developed an initial algorithm for the defensive robots.
\end{itemize}

\textbf{Positioning (and what is novel here).} Task allocation in multi-robot systems has been widely studied, including market/bidding-style assignments and distributed coordination under limited communication. Separately, formation/consensus control is a standard tool for maintaining team structure using local neighbor interactions. Our contribution is not a new core algorithmic primitive; rather, we combine (i) local formation maintenance, (ii) a simple distributed ``bidding'' proxy for primary fielder selection, and (iii) a relay-vs-direct delivery decision into one end-to-end pipeline tailored to a baseball-like infield scenario, with an implementation that makes the architecture and evaluation protocol explicit. This integration and evaluation in an infield-defense-inspired scenario is the primary novelty claim, and we compare against simple baselines in Section~VII.

\subsection{Related Work (brief)}
Multi-robot task allocation has been formalized and categorized in prior work (e.g., auction/market-based methods)~\cite{gerkey2004}. Consensus and graph-theoretic methods are commonly used to maintain coordination and distributed agreement in multi-agent teams~\cite{mesbahi2010,olfati2006}. Our approach uses these ideas in a deliberately simple form to prioritize a complete, testable end-to-end system in simulation: local formation control plus a distributed selection rule for interception and a one-hop relay decision for delivery.

\section{System Model}

Consider a team of $N$ robots operating in a two-dimensional environment. The position of robot $i$ is defined as:

\begin{equation}
p_i = [x_i, y_i]^T \in \mathbb{R}^2 .
\end{equation}

The position of the ball at time $t$ is:

\begin{equation}
p_b(t) = [x_b(t), y_b(t)]^T \in \mathbb{R}^2 .
\end{equation}

Robot motion follows a kinematic model:

\begin{equation}
\dot{p}_i = v_i ,
\end{equation}
with velocity constrained by $\|v_i\| \le v_{\max}$ when a speed limit is enforced.

The objective is to minimize total task completion time:

\begin{equation}
T_{\mathrm{total}} = T_{\mathrm{intercept}} + T_{\mathrm{delivery}} .
\end{equation}

\section{Implemented Approach}

\subsection{Ball Trajectory Model}

The ground ball trajectory is approximated as:

\begin{equation}
p_b(t) = p_0 + v_0 t + \frac{1}{2} a t^2 , \qquad t \ge 0 ,
\end{equation}
where \(p_0 \in \mathbb{R}^2\) is the ball position at the start of the segment, \(v_0 \in \mathbb{R}^2\) is the initial planar velocity, and \(a \in \mathbb{R}^2\) is the (constant) planar acceleration used to approximate rolling friction and small bounces. Here \(t\) is elapsed time since the segment began. This allows robots to query predicted ball positions \(p_b(t)\) from the model.

\subsection{Formation Control}

Let $\mathcal{N}_i$ denote the set of neighbor indices of robot $i$ on the communication graph. We implement a decentralized control law:
\begin{equation}
u_i = -k \sum_{j \in \mathcal{N}_i} (p_i - p_j) ,
\label{eq:formation-law}
\end{equation}
with gain $k > 0$. This is a Laplacian-type (consensus) coupling: it coordinates relative positions using only local information. It does not, by itself, enforce fixed pairwise distances; spacing behavior depends on the graph, gains, and initial conditions.

\subsection{Interception Model}

At a common evaluation time $t$, each robot $i$ uses a straight-line travel-time proxy
\begin{equation}
\widehat{T}_i(t) = \frac{\|p_i - p_b(t)\|}{v_{\max}} .
\label{eq:Ti-hat}
\end{equation}
\noindent\textbf{Modeling assumption (and justification).} We treat robots as holonomic point agents with bounded speed \(v_{\max}\) and neglect acceleration limits and turning dynamics. Under this model, the minimum time to reach a \emph{fixed} point \(q\) is \(\|p_i-q\|/v_{\max}\). For a moving ball, using \(\|p_i-p_b(t)\|/v_{\max}\) at a shared evaluation time \(t\) becomes a lightweight heuristic ``bid'' that is computable locally and is consistent with the simulator's discrete-time update loop. In the final report we will compare this proxy against an alternative that searches over a short horizon for the best intercept time \(t'\) (minimizing \(\|p_i-p_b(t')\|/v_{\max}+t'\)) to quantify the approximation impact.
\noindent The selected fielder is
\begin{equation}
i^\ast = \arg\min_i \widehat{T}_i(t) .
\label{eq:istar}
\end{equation}

\subsection{Distributed Coordination}

The system implements:
\begin{itemize}
    \item Interception bidding and selection
    \item Role assignment (primary and supporting robots)
    \item State transitions between formation, interception, and delivery
\end{itemize}

\subsection{Delivery Optimization}

After interception, the system evaluates direct travel versus a one-hop relay using straight-line legs at speed $v_{\max}$. Let $p_{\mathrm{base}}$ denote the target base position. Let $\tau_{\mathrm{hp}} \ge 0$ denote any fixed handling/pass delay at the relay (instantaneous handoff corresponds to $\tau_{\mathrm{hp}} = 0$):
\begin{equation}
C_{\mathrm{direct}} = \frac{\|p_{i^\ast} - p_{\mathrm{base}}\|}{v_{\max}} ,
\end{equation}
\begin{equation}
C_{\mathrm{relay}}(j) = \frac{\|p_{i^\ast} - p_j\|}{v_{\max}} + \frac{\|p_j - p_{\mathrm{base}}\|}{v_{\max}} + \tau_{\mathrm{hp}} .
\end{equation}
\noindent\textbf{Modeling assumption (and justification).} The delivery model assumes the same holonomic, speed-bounded travel-time proxy used for interception. The relay cost adds a constant handoff delay \(\tau_{\mathrm{hp}}\) to represent capture/transfer latency; in simulation this can be set to \(0\) (idealized) or tuned to study when relays stop being beneficial. This yields a transparent criterion for whether ``throwing''/handoff is worth it, and it is directly measurable in repeated trials.
\noindent The system selects $\min\bigl\{ C_{\mathrm{direct}},\; \min_{j \ne i^\ast} C_{\mathrm{relay}}(j) \bigr\}$ (and the corresponding action).

\subsection{Technical Implementation}

\begin{itemize}
    \item Language: Python
    \item Simulation: Custom 2D environment
    \item Modules:
    \begin{itemize}
        \item Trajectory prediction
        \item Interception solver
        \item Communication layer (broadcast)
        \item Role assignment logic
    \end{itemize}
\end{itemize}

\section{System Architecture}

\noindent\textbf{Distributed algorithm vs. centralized simulator harness.} Our \emph{algorithmic} claim is distributed: each robot can compute its own bid \(\widehat{T}_i(t)\) and (given a communication primitive to share bids) agree on \(i^\ast\), and each robot can evaluate relay candidates using locally available states. Our current \emph{implementation} is a centralized Python simulator that executes a synchronous ``round'' each timestep (to update dynamics, sample noise, and render plots). This centralized harness collects the robots' per-agent outputs for convenience, but the computations labeled \texttt{formation} and \texttt{coordination} are written in a way that is per-robot computable given neighbor/broadcast messages. In the final version, we will make this separation explicit by (i) treating message passing as an interface and (ii) reporting which variables are local (\(p_i\), neighbor states) versus global (ball state, for this simulator).

\subsection{Computation and Information Partition}
To remove ambiguity about ``local'' versus ``global'' logic, we define the following partition:
\begin{itemize}
    \item \textbf{Local robot state (robot \(i\)):} \(p_i\) (and \(v_i\) if modeled), plus its current mode (formation / pursuit / delivery).
    \item \textbf{Neighbor information:} \(p_j\) for \(j \in \mathcal{N}_i\) (used in \eqref{eq:formation-law}).
    \item \textbf{Broadcast information (shared):} a detection timestamp and a shared estimate of ball parameters \((p_0,v_0,a,t_0)\) for \(p_b(t)\). In real hardware, this could come from a sensor robot or a shared perception service; in our simulator it is provided directly by \texttt{ball}.
    \item \textbf{Distributed coordination messages:} scalar bids \(\widehat{T}_i(t)\) and optional relay costs \(C_{\mathrm{relay}}(j)\). An ``all-to-all'' broadcast is assumed for the current simulation; restricting this to a neighborhood graph is a possible extension.
\end{itemize}
\noindent Under this partition, the only centralized component is the \emph{simulation loop} that advances time and applies dynamics; the decision rules themselves are defined per-agent and depend only on local/communicated quantities.

Figure~\ref{fig:arch} summarizes the software structure. Robots are treated as point masses in the plane. Each simulation step updates either defensive formation or pursuit behavior depending on game phase.

\begin{figure}[htbp]
\centering
\fbox{\parbox{0.88\columnwidth}{\centering
\vspace{0.15in}
\textbf{infield\_defense} package (Python)\\[0.5em]
\texttt{config} $\rightarrow$ parameters (field, $v_{\max}$, $k$, $\Delta t$)\\[0.25em]
\texttt{ball} $\rightarrow$ $p_b(t)$ prediction\\[0.25em]
\texttt{formation} $\rightarrow$ neighbor graph + $u_i$\\[0.25em]
\texttt{coordination} $\rightarrow$ $\widehat{T}_i$, $i^\ast$, $C_{\mathrm{direct}}$, $C_{\mathrm{relay}}$\\[0.25em]
\texttt{simulation} $\rightarrow$ state, Euler integration, ball launch\\[0.25em]
\texttt{cli} $\rightarrow$ Matplotlib animation
\vspace{0.15in}
}}
\caption{Modular architecture of the simulation. Replace the framed box with a TikZ or PowerPoint figure if required.}
\label{fig:arch}
\end{figure}

\noindent\textbf{Data flow.} \texttt{simulation} holds robot positions and velocities and the active \texttt{BallState} $(p_0, v_0, a, t_0)$. On ball detection, \texttt{coordination.select\_primary\_fielder} sets $i^*$. During pursuit, the demo steers $i^*$ toward the current $p_b(t)$; supporting robots continue formation on the induced subgraph. Delivery optimization is implemented but not yet wired to post-capture motion.

\section{Experimental Design}

To evaluate the performance of the proposed system, we design a series of simulation experiments.

\subsection{Simulation Setup}

The environment consists of $N$ robots initialized in a defensive formation. The ball is launched with randomized initial velocity and direction. Each simulation will run until interception and delivery are completed.

\subsection{Variables}

We vary the following parameters:
\begin{itemize}
    \item Number of robots ($N = 4, 6, 8$)
    \item Maximum robot velocity ($v_{\max}$)
    \item Ball initial speed and direction
    \item Communication delays
\end{itemize}

\subsection{Evaluation Metrics}

We measure:
\begin{itemize}
    \item Interception time ($T_{\text{intercept}}$)
    \item Delivery time ($T_{\text{delivery}}$)
    \item Total task completion time ($T_{\text{total}}$)
    \item Success rate (successful fielding and delivery)
\end{itemize}

\subsection{Evaluation Protocol (repeated trials and baselines)}
For each parameter setting, we will run \(K\) randomized trials (we target \(K=50\) or \(100\)) with controlled random seeds and report mean \(\pm\) standard deviation as well as median and 90th percentile for timing metrics. We will include at least two baselines:
\begin{itemize}
    \item \textbf{Nearest-direct baseline:} choose \(i^\ast\) by Euclidean distance to the ball at detection; deliver directly (no relay).
    \item \textbf{Ours w/o relay:} our bidding-based \(i^\ast\), but force direct delivery. This isolates the relay decision benefit.
\end{itemize}

\subsection{Expected Outcomes}

We think that increasing the number of robots improves interception success, but could introduce coordination overhead. In addition, relay-based delivery is expected to outperform direct delivery in long-distance scenarios.

\section{Preliminary Results}

We have verified the implementation qualitatively via the Matplotlib demo: robots maintain an arc-shaped defensive layout under formation control; after a scripted ground ball launch, the primary fielder index matches $\arg\min_i \widehat{T}_i(t)$ under the proxy~\eqref{eq:Ti-hat}. Quantitative evaluation is in progress; the next deliverable is a table of repeated-trial results with baseline comparisons.

\subsection{Quantitative Results Template (to be filled)}
Table~\ref{tab:results-template} is the intended format for reporting repeated trials. We will also include plots of distributions (histograms/ECDFs) for \(T_{\mathrm{intercept}}\) and \(T_{\mathrm{total}}\) and a success-rate curve versus \(v_{\max}\).

\begin{table}[htbp]
\centering
\caption{Preliminary quantitative results template (\(K\) trials per setting).}
\label{tab:results-template}
\begin{tabular}{@{}lcccc@{}}
\toprule
Method & \(N\) & \(v_{\max}\) & \(T_{\mathrm{intercept}}\) (mean\(\pm\)std) & Success (\%) \\
\midrule
Nearest-direct baseline &  &  &  &  \\
Ours (no relay) &  &  &  &  \\
Ours (relay enabled) &  &  &  &  \\
\bottomrule
\end{tabular}
\end{table}

\subsection{Current Evidence}
As an existence proof of the full simulation loop, Figure~\ref{fig:placeholder} shows the running system. This figure will be complemented by quantitative tables/plots as described above.

\begin{figure}
    \centering
    \includegraphics[width=1\linewidth]{Simulation Screenshot.jpg}
    \caption{Simulation screenshot (qualitative verification of the pipeline).}
    \label{fig:placeholder}
\end{figure}

\section{Plan and Milestones}

\subsection{Remaining Tasks}

\begin{itemize}
    \item Improve communication/coordination between robots
    \item Finish algorithm for ball interception
    \item Run repeated-trial evaluation and generate quantitative tables/plots (timings and success rates) with baseline comparisons
    \item Integrate delivery (direct vs relay) into post-capture motion and report when relay helps as a function of distance and \(\tau_{\mathrm{hp}}\)
\end{itemize}

\subsection{Evaluation Plan}

\textbf{Scenarios:}
\begin{itemize}
    \item Varying ball speeds
    \item Different robot counts
    \item Communication delays
    \item Formation disturbances
\end{itemize}

\textbf{Metrics:}
\begin{itemize}
    \item Efficiency metrics (completion time, success rate)
    \item Formation stability
    \item Communication quality
\end{itemize}

\subsection{Success Criteria}
We will consider the core system complete if:
\begin{itemize}
    \item For at least one representative setting (e.g., \(N \in \{4,6\}\) and two values of \(v_{\max}\)), we can report \(K \ge 50\) trial results with mean\(\pm\)std for \(T_{\mathrm{intercept}}\), \(T_{\mathrm{delivery}}\), \(T_{\mathrm{total}}\), plus success rate.
    \item Our method matches or improves \(T_{\mathrm{total}}\) compared to the nearest-direct baseline in at least one regime (e.g., longer delivery distance or higher \(N\)), and we can explain when it fails.
    \item Plots/tables and figure captions are sufficient for an academic reader to verify claims without relying on screenshots alone.
\end{itemize}

\section{Conclusion}

Our project presents a decentralized framework for multi-robot coordination inspired by baseball infield defense. By combining formation control, interception optimization, and relay-based decision-making, the system demonstrates how distributed algorithms can effectively handle dynamic task allocation.

\begin{thebibliography}{00}

\bibitem{mesbahi2010} M. Mesbahi and M. Egerstedt, \textit{Graph Theoretic Methods in Multiagent Networks}, Princeton University Press, 2010.

\bibitem{gerkey2004} B. P. Gerkey and M. J. Mataric, ``A formal analysis and taxonomy of task allocation in multi-robot systems,'' \textit{International Journal of Robotics Research}, 2004.

\bibitem{olfati2006} R. Olfati-Saber, ``Flocking for multi-agent dynamic systems,'' \textit{IEEE Transactions on Automatic Control}, 2006.

\end{thebibliography}

\end{document}
